% =============================================================================
% Technical report: the constrained dual-arm task-space BiRRT (CDFM planner)
% as implemented in husky_assembly_tamp.motion_planner.
%
% Build:  pdflatex cdfm_task_space_birrt_report.tex   (run twice for refs)
% =============================================================================
\documentclass[11pt,a4paper]{article}

% --- page + fonts ------------------------------------------------------------
\usepackage[margin=2.2cm]{geometry}
\usepackage[T1]{fontenc}
\usepackage[utf8]{inputenc}
\usepackage{lmodern}

% --- maths -------------------------------------------------------------------
\usepackage{amsmath,amssymb}

% --- pseudocode ---------------------------------------------------------------
\usepackage{algorithm}
\usepackage{algpseudocode}
\usepackage{caption}

% --- tables, colours, links ---------------------------------------------------
\usepackage{booktabs}
\usepackage{array}
\usepackage{xcolor}
\usepackage{enumitem}
\usepackage[colorlinks=true,linkcolor=blue!60!black,urlcolor=blue!60!black]{hyperref}

% =============================================================================
% Macros
% =============================================================================

% Code identifiers (underscores are typed as-is thanks to \detokenize).
\newcommand{\code}[1]{\texttt{\detokenize{#1}}}

% A line of pseudocode that DIFFERS from the standard BiRRT is printed in red
% and gets a [D-n] tag at the right margin, pointing to the difference list
% in Section 5.
\definecolor{diff}{RGB}{176,32,32}
\newcommand{\D}[1]{\textcolor{diff}{#1}}
\newcommand{\dtag}[1]{\hfill\textcolor{diff}{\scriptsize\textsf{[#1]}}}

% Pseudocode comments in grey, to the right.
\algrenewcommand{\algorithmiccomment}[1]{\hfill\textcolor{gray}{\footnotesize$\triangleright$ #1}}
\newcommand{\Continue}{\State \textbf{continue}}
\newcommand{\Break}{\State \textbf{break}}

% Non-floating algorithm block that may break across pages.
\newenvironment{algo}[2]{%
  \par\medskip\noindent\rule{\linewidth}{0.8pt}\par
  \captionof{algorithm}{#1}\label{#2}%
  \vspace{-4pt}\noindent\rule{\linewidth}{0.4pt}\par\vspace{-2pt}
  \begin{algorithmic}[1]}%
 {\end{algorithmic}\vspace{-2pt}\noindent\rule{\linewidth}{0.8pt}\par\medskip}

% Maths shorthands.
\newcommand{\R}{\mathbb{R}}
\newcommand{\SE}{\mathrm{SE}(3)}
\newcommand{\norm}[1]{\left\lVert #1 \right\rVert}
\newcommand{\Tree}{\mathcal{T}}
\newcommand{\fn}[1]{\textsc{#1}}

\title{The Constrained Dual-Arm Task-Space BiRRT\\
       \large A technical report on the current CDFM (M1) planner implementation}
\author{husky\_assembly\_tamp \texttt{motion\_planner} --- written for Yijiang Huang}
\date{29 September 2026}

\begin{document}
\maketitle

\begin{abstract}
\noindent
This report explains the bidirectional task-space RRT that plans the M1
``constrained dual-arm free movement'' (CDFM): both arms of the Husky dual-UR5e
hold one bar rigidly and carry it from a home (loading) pose to the approach
pose in front of the structure. The search runs on the \emph{bar pose} in
$\SE$, and every tree node is lifted to a 12-joint configuration by analytical
inverse kinematics as the tree grows. The pseudocode below is written directly
from the code in \code{dual_arm_task_space_rrt/core.py} and
\code{api.py}; lines that differ from the textbook RRT-Connect of
Kuffner and LaValle are printed in \D{red} and tagged \D{\textsf{[D-n]}},
with the list of differences collected in Section~\ref{sec:diff}.
\end{abstract}

\tableofcontents

% =============================================================================
\section{What is being planned}
\label{sec:problem}
% =============================================================================

\paragraph{The movement.}
An M1 movement starts with the bar already gripped by both tools and ends with
the bar at its approach pose next to the partially built structure. Nothing is
grasped or released along the way, so the two tool frames keep a fixed
relation to the bar for the whole motion. The robot base does not move.

\paragraph{Variables.}
\begin{itemize}[nosep]
  \item $q = (q_L, q_R) \in \R^{12}$: the six left-arm joints followed by the
        six right-arm joints (canonical UR order).
  \item $T \in \SE$: the world pose of the bar. In code a pose is a PyBullet
        pair \code{(pos, quat_xyzw)}.
  \item $G_L, G_R \in \SE$: the two \emph{grasp transforms}, bar-from-tool0.
        They are read once from the goal (or start) configuration and never
        change. The world tool poses implied by a bar pose are
        $T\,G_L$ and $T\,G_R$.
  \item $\mathrm{FK}_L(q_L), \mathrm{FK}_R(q_R)$: forward kinematics of each
        arm to its \code{tool0} link.
\end{itemize}

\paragraph{The constraint.}
A configuration is admissible only when both tools sit exactly on the bar:
\begin{equation}
  \mathcal{M} \;=\; \bigl\{\, q \in \R^{12} \;:\;
    \exists\, T \in \SE \text{ with } \mathrm{FK}_L(q_L) = T G_L
    \text{ and } \mathrm{FK}_R(q_R) = T G_R \,\bigr\}.
  \label{eq:manifold}
\end{equation}
$\mathcal{M}$ is a six-dimensional surface inside the twelve-dimensional joint
space; a random joint sample lands on it with probability zero. That is why the
planner does not sample joints at all. It samples the bar pose $T$ (six
numbers), and \emph{lifts} each $T$ to a $q \in \mathcal{M}$ by solving
inverse kinematics for both arms. The only freedom hidden inside that lift is
the choice of IK branch (a UR arm has up to eight for one tool pose), and much
of the machinery described below exists to keep that choice consistent along a
path.

\paragraph{What ``solved'' means.}
The output is a list of bar poses $T_0,\dots,T_n$ with $T_0 = T_{\text{start}}$,
$T_n = T_{\text{goal}}$, and the matching joint path $q_0,\dots,q_n$ such that
\begin{enumerate}[nosep,label=(\roman*)]
  \item every $q_i$ holds the bar at $T_i$ (Eq.~\ref{eq:manifold}, verified by
        FK to $10^{-4}$),
  \item consecutive poses are at most one step
        ($\delta_p$ metres, $\delta_r$ radians) apart,
  \item consecutive configurations are at most $\tau$ apart per joint,
        $\norm{q_{i+1}-q_i}_\infty \le \tau$ with $\tau = 10^\circ$
        (the \emph{continuity gate}), and
  \item every $q_i$ is collision-free for the whole cell (robot, tools,
        attached bar, structure, obstacles).
\end{enumerate}

% =============================================================================
\section{Where the code lives}
\label{sec:where}
% =============================================================================

\begin{center}
\footnotesize
\begin{tabular}{@{}llp{7.2cm}@{}}
\toprule
\textbf{Function} & \textbf{File} & \textbf{Role} \\
\midrule
\code{plan_constrained_dual_arm} & \code{api.py} & public entry point; endpoints, stage switches, calls the search, then the smoother \\
\code{_derive_constrained_start_for_plan} & \code{api.py} & goal IK, goal/home branch pairing, backward-tracked start, direct-path shortcut \\
\code{_ssik_pair_goal_branch_with_home} & \code{api.py} & pick the goal IK branch pair that a home pose can reach \\
\code{plan_pose_birrt} & \code{core.py} & \textbf{the bidirectional search} (Alg.~\ref{alg:birrt}) \\
\code{extend_toward} & \code{core.py} & greedy walk of one tree toward a pose (Alg.~\ref{alg:extend}) \\
\code{solve_dual_arm_pose_ik} & \code{core.py} & lift one bar pose to $q$ (Alg.~\ref{alg:ik}) \\
\code{joint_step_exceeds_threshold} & \code{core.py} & the continuity gate \\
\code{pose_to_feature_vec}, \code{nearest_node} & \code{core.py} & the metric and the nearest-neighbour scan \\
\code{sample_pose} & \code{core.py} & the random bar-pose sampler \\
\code{build_cfab_pose_collision_fn} & \code{core.py} & bar-only collision pre-check (CC.4) \\
\code{derive_constrained_start_tracked} & \code{core.py} & start derivation by tracking goal$\to$home (Alg.~\ref{alg:tracked}) \\
\code{smooth_dual_arm_pose_path} & \code{smooth.py} & shortcut smoother in pose space (Alg.~\ref{alg:smooth}) \\
\code{ssik_arm_branches}, \code{rebranch_toward_seed} & \code{ssik_ik.py} & analytical IK branches, $2\pi$ re-branching (Alg.~\ref{alg:branches}) \\
\code{ssik_inprocess.solve} & \code{keyframe/} & loads the generated per-arm solver module and forwards the call \\
artifact \code{solve}, \code{_solve_algebraic} & \code{asset/ssik/} & the generated per-arm ssik solver: candidates, FK gate, dedup, limits (Algs.~\ref{alg:artifact}, \ref{alg:rr}) \\
\bottomrule
\end{tabular}
\end{center}

\noindent
All paths are relative to \path{husky_assembly_tamp/motion_planner/} unless
they start with \path{keyframe/} or \path{asset/}, which are siblings of
\path{motion_planner/} inside the package; the RRT files sit in the
\path{dual_arm_task_space_rrt/} sub-package. Line numbers
quoted in Section~\ref{sec:reading} refer to the tree at the time of writing.
An earlier, Markdown-format review of the same code is in
\path{docs/plan_constrained_dual_arm.md}; the benchmark that measures it is
described in \path{benchmark/README.md} and \path{benchmark/analysis_night1.md}.

% =============================================================================
\section{The textbook baseline: RRT-Connect}
\label{sec:baseline}
% =============================================================================

For reference, Algorithm~\ref{alg:std} is RRT-Connect as introduced by Kuffner
and LaValle (ICRA 2000), which is what people usually mean by ``BiRRT''. It
plans in the configuration space $\mathcal{C}$ with a single validity test
(\fn{CollisionFree}) and a fixed step $\varepsilon$.

\begin{algo}{Standard RRT-Connect (Kuffner \& LaValle 2000)}{alg:std}
\Function{RRT-Connect}{$q_{\text{init}}, q_{\text{goal}}$}
  \State $\Tree_a \gets \{q_{\text{init}}\}$;\quad $\Tree_b \gets \{q_{\text{goal}}\}$
  \For{$k = 1 \dots K$}
    \State $q_{\text{rand}} \gets$ \fn{RandomConfig}() \Comment{uniform in $\mathcal{C}$}
    \If{\fn{Extend}($\Tree_a$, $q_{\text{rand}}$) $\neq$ \textsc{Trapped}}
      \If{\fn{Connect}($\Tree_b$, $q_{\text{new}}$) $=$ \textsc{Reached}}
        \State \Return \fn{Path}($\Tree_a$, $\Tree_b$) \Comment{concatenate the two branches at $q_{\text{new}}$}
      \EndIf
    \EndIf
    \State \fn{Swap}($\Tree_a$, $\Tree_b$) \Comment{strict alternation}
  \EndFor
  \State \Return \textsc{Failure}
\EndFunction
\Statex
\Function{Extend}{$\Tree$, $q$}
  \State $q_{\text{near}} \gets$ \fn{Nearest}($\Tree$, $q$) \Comment{Euclidean in $\mathcal{C}$}
  \State $q_{\text{new}} \gets q_{\text{near}} + \varepsilon\,\dfrac{q - q_{\text{near}}}{\norm{q - q_{\text{near}}}}$ \Comment{one fixed step}
  \If{\fn{CollisionFree}($q_{\text{near}}, q_{\text{new}}$)}
    \State add vertex $q_{\text{new}}$ and edge $(q_{\text{near}}, q_{\text{new}})$ to $\Tree$
    \State \Return \textsc{Reached} if $q_{\text{new}} = q$ else \textsc{Advanced}
  \EndIf
  \State \Return \textsc{Trapped}
\EndFunction
\Statex
\Function{Connect}{$\Tree$, $q$}
  \Repeat
    \State $S \gets$ \fn{Extend}($\Tree$, $q$)
  \Until{$S \neq$ \textsc{Advanced}} \Comment{keep stepping until reached or blocked}
  \State \Return $S$
\EndFunction
\end{algo}

\noindent
Two properties of the baseline matter for the comparison:
\begin{itemize}[nosep]
  \item A node \emph{is} a configuration; there is nothing to compute when a
        node is added, and concatenating the two trees at the meeting node is
        trivially a valid path.
  \item Validity is one boolean (collision). There is no notion of a step
        being \emph{kinematically} inconsistent with the previous one.
\end{itemize}

% =============================================================================
\section{The implemented search}
\label{sec:impl}
% =============================================================================

\subsection{Building blocks}

\paragraph{Pose interpolation.}
$\fn{Interp}(T_a, T_b;\ \delta_p, \delta_r)$ (PyBullet-planning
\code{interpolate_poses}) returns
$N = \lceil \max(\norm{p_b - p_a}/\delta_p,\ \theta(R_a,R_b)/\delta_r) \rceil$
intermediate poses with linear position and slerp orientation, followed by
$T_b$ itself. So the last element of the walk is \emph{exactly} the target
pose; this is what makes ``the trees meet'' an equality rather than a
tolerance test in the bidirectional planner.

\paragraph{Feature-point metric.}
The bar's eight bounding-box corners $c_1,\dots,c_8$ (bar-local; the box is
$0.03 \times 0.03 \times 1.0$ m by default) are pushed through a pose to give
a 24-vector
\begin{equation}
  \phi(T) = \bigl[\,T c_1;\ T c_2;\ \dots;\ T c_8\,\bigr] \in \R^{24},
  \qquad
  d(T_1, T_2) = \norm{\phi(T_1) - \phi(T_2)}_2 .
  \label{eq:feature}
\end{equation}
This single number blends translation and rotation in a way that is scaled by
the physical size of the bar (rotating a long bar moves its corners a lot),
which avoids having to pick a weight between metres and radians. $\phi$ is
cached per node so nearest-neighbour search is a linear scan over cached
24-vectors.

\paragraph{Continuity gate.}
\begin{equation}
  \fn{Jump}(q', q) \;\equiv\; \norm{q' - q}_\infty > \tau,
  \qquad \tau = \code{DEFAULT_JOINT_CONTINUITY_THRESHOLD_RAD} = 10^\circ .
  \label{eq:gate}
\end{equation}
The comparison is on \emph{raw} joint values (no wrapping to $\pm\pi$; the
unwrapping helper is left in the source but disabled). A step that fails this
gate is interpreted as the IK having flipped to a different solution branch,
and is discarded.

\paragraph{Two collision predicates.}
\begin{itemize}[nosep]
  \item \code{collision_fn}$(T)$: bar-only. The attached bar is moved to the
        sampled world pose by rewriting its attachment frame, and only
        compas\_fab's step CC.4 (attached rigid body vs.\ environment) runs.
        Cheap pre-reject, used in every stage.
  \item \code{joint_collision_fn}$(q)$: full cell. Clones the start state,
        writes the twelve joints, runs the complete compas\_fab check (self,
        tools, attached bar, environment, ACM). Stage 3 only. Measured cost
        $\approx 13$ ms per call (Section~\ref{sec:bench}).
\end{itemize}

\subsection{Lifting a bar pose to joints}

\begin{algo}{\textsc{SolveDualArmPoseIK} --- \texttt{solve\_dual\_arm\_pose\_ik} (ssik backend)}{alg:ik}
\Function{SolveDualArmPoseIK}{$T$, $q_{\text{seed}}$}
  \For{$\text{arm} \in \{L, R\}$}
    \State $X \gets T\,G_{\text{arm}}$ \Comment{world tool0 target}
    \State $B \gets$ \code{ssik_arm_branches}(arm, $X$, $q_{\text{seed},\text{arm}}$, \code{allow_rescue=False})
    \Statex \hspace{\algorithmicindent}\hspace{\algorithmicindent}
      \textcolor{gray}{\footnotesize every analytical branch inside joint limits, each shifted by $k\cdot 2\pi$ per joint toward the seed when that stays in-limit, sorted by $\norm{\cdot - q_{\text{seed},\text{arm}}}_\infty$ (Alg.~\ref{alg:branches}, Section~\ref{sec:ssik})}
    \If{$B = \varnothing$} \Return \textsc{None} \EndIf \Comment{tool pose unreachable for this arm}
    \State $q_{\text{arm}} \gets B[0]$ \Comment{nearest branch to the seed}
  \EndFor
  \State $q \gets (q_L, q_R)$
  \If{\textbf{not} \code{validate_dual_arm_bar_pose}($q$, $T$)} \Return \textsc{None} \EndIf
  \Statex \hspace{\algorithmicindent}\textcolor{gray}{\footnotesize PyBullet FK re-check: both tools within $10^{-4}$ of $T G_L$, $T G_R$, and the two implied bar poses agree}
  \State \Return $q$
\EndFunction
\end{algo}

\noindent
Two remarks. First, the two arms are independent 6R chains, so the lift is two
separate analytical solves; nothing about the second arm depends on the first.
Second, ``nearest branch to the seed'' is the \emph{only} mechanism that keeps
a walk on one branch sheet. There is no memory of which branch was chosen
earlier beyond the previous configuration; the continuity gate
(Eq.~\ref{eq:gate}) is the safety net that catches the cases where nearest
still means far. What ``branch'' means, and how the list $B$ is produced, is
the subject of the next subsection.

\subsection{Inside \texttt{ssik\_arm\_branches}: how the IK branches are produced}
\label{sec:ssik}

Every IK call in the planner goes through one function,
\code{ssik_ik.ssik_arm_branches}. It is a thin adapter over the analytical
solver library \emph{ssik} (v3.2.0 in the planner's venv), and the stack it
sits on has four layers:

\begin{center}
\footnotesize
\begin{tabular}{@{}p{1.3cm} p{5.0cm} p{9.2cm}@{}}
\toprule
layer & where & what it does \\
\midrule
adapter & \path{motion_planner/ssik_ik.py} & world $\to$ arm-base frame, $2\pi$ re-branching toward the seed, raw-distance sort (Alg.~\ref{alg:branches}) \\
\addlinespace
loader & \path{keyframe/ssik_inprocess.py} & imports the generated per-arm module once and forwards the call \\
\addlinespace
artifact & \path{asset/ssik/left_ur_arm_ik.py}, \path{asset/ssik/right_ur_arm_ik.py} & the generated solver: algebraic candidates, FK check, dedup, limits, seed sort (Alg.~\ref{alg:artifact}) \\
\addlinespace
library & \path{ssik/solvers/ikgeo/}, \path{ssik/postprocess.py}, \path{ssik/refinement/} (installed package) & Raghavan--Roth elimination, limit filters, perturbation rescue (Algs.~\ref{alg:rr}, \ref{alg:rescue}) \\
\bottomrule
\end{tabular}
\end{center}

\paragraph{Why it is not the textbook UR solver.}
The artifact modules were generated once by \code{ssik build} from the
\emph{calibrated} URDF, with each arm rooted at its own
\code{*_ur_arm_base_link} and ending at \code{*_ur_arm_tool0}. The familiar
closed-form UR inverse kinematics (eight postures: shoulder left/right, elbow
up/down, wrist flipped/not) relies on the nominal geometry -- axes 2, 3, 4
exactly parallel, axes 4, 5, 6 pairwise perpendicular and meeting. Calibration
breaks that. In the left-arm artifact the second joint axis is
$(2.7\times10^{-6},\ 0.9999998,\ -5.9\times10^{-4})$ and the fourth is
$(-1.2\times10^{-3},\ 0.99997,\ -7.9\times10^{-3})$: parallel to within about
half a degree, but not parallel. ssik's dispatcher therefore found ``no tier-0
(Pieper-class) match'' and emitted its \emph{tier-2 general 6R} solver,
\code{ikgeo.general_6r}: the Raghavan--Roth elimination with the
Manocha--Canny eigenvalue formulation. That solver makes no structural
assumption about the chain and can return up to 16 real solutions for a
general 6R arm; for a chain this close to the UR geometry one expects the
usual eight postures among them, but the code never assumes that count. The
artifact records \code{EXPECTED_MS_MEDIAN = 5.0} ms per solve.

\paragraph{Joint limits that the solver enforces.}
The artifact carries the URDF limits, in radians: shoulder pan
$\pm 2\pi$, shoulder lift $\pm 4.712$ ($\pm 270^\circ$), elbow $\pm\pi$, and
each wrist $\pm 4.189$ ($\pm 240^\circ$). Four of the six ranges are wider
than one full turn, so the same physical posture can be represented by two
(pan: up to three) different in-limit joint vectors that differ by $2\pi$. The
elbow range is exactly one turn, so the elbow never has an in-limit twin. This
is the root of the $2\pi$ issue handled by the adapter.

\paragraph{Vocabulary.}
A \emph{branch} is one geometrically distinct solution: two joint vectors that
agree modulo $2\pi$ on every joint are the same branch. A \emph{$2\pi$ twin}
is a different representative of the same branch. ssik returns one
representative per branch; the adapter's job is to pick the representative
that sits nearest to the seed in \emph{raw} joint values, because that is what
the planner's continuity gate measures.

\begin{algo}{\textsc{SsikArmBranches} --- \texttt{ssik\_ik.ssik\_arm\_branches} and \texttt{rebranch\_toward\_seed}}{alg:branches}
\Function{SsikArmBranches}{arm, $X_{\text{world}}$, $q_{\text{seed}}$, \code{max_solutions}, \code{allow_rescue}}
  \State \textsc{VerifyJointOrder}(robot) \Comment{once per robot: artifact joint order = URDF chain order}
  \State $W \gets$ world pose of \code{<arm>_ur_arm_base_link}, read live from PyBullet
  \State $T \gets W^{-1}\, X_{\text{world}}$ \Comment{ssik expects tool0 in the arm-base frame}
  \State sols $\gets$ \textsc{ArtifactSolve}($T$; \code{max_solutions}, $q_{\text{seed}}$, \code{allow_rescue}) \Comment{Alg.~\ref{alg:artifact}}
  \State $B \gets [\,q \text{ for } (q, \text{residual}) \in \text{sols}\,]$
  \If{$q_{\text{seed}} \neq$ \textsc{None}}
    \State \D{$B \gets [\,\textsc{Rebranch}(q, q_{\text{seed}}, \text{limits}) \text{ for } q \in B\,]$}
    \State \D{sort $B$ by $\norm{q - q_{\text{seed}}}_\infty$ ascending} \Comment{raw distance, the metric of Eq.~\ref{eq:gate}}
  \EndIf
  \State \Return $B$ \Comment{empty = unreachable, or every branch outside the limits}
\EndFunction
\Statex
\Function{Rebranch}{$q$, seed, limits}
  \For{$i = 1 \dots 6$}
    \State $k \gets \operatorname{round}\bigl((\text{seed}_i - q_i)/2\pi\bigr)$ \Comment{whole turns that bring $q_i$ onto the seed's turn}
    \If{$\text{lo}_i \le q_i + 2\pi k \le \text{hi}_i$} $q_i \gets q_i + 2\pi k$ \EndIf \Comment{only when the twin is inside the limits}
  \EndFor
  \State \Return $q$
\EndFunction
\end{algo}

\noindent
The two red lines are the whole point of the adapter. ssik already sorts by
distance to the seed, but with a \emph{wrap-to-$\pi$} metric (its
\code{seed_metric="wrap_linf"}), under which a joint at $3.0$ rad and its twin
at $3.0 - 2\pi = -3.28$ rad are at distance zero. The planner's continuity gate
takes the raw difference, $6.28$ rad, and would reject that step as a branch
flip. \textsc{Rebranch} moves every joint onto the seed's turn when the limits
allow it, and the re-sort makes the first element the branch that is nearest
in the sense the planner will actually test.

\begin{algo}{\textsc{ArtifactSolve} --- the generated \texttt{solve()} in \texttt{asset/ssik/<arm>\_ur\_arm\_ik.py}}{alg:artifact}
\Function{ArtifactSolve}{$T$; \code{max_solutions}, $q_{\text{seed}}$, \code{allow_rescue}}
  \State candidates $\gets$ \textsc{SolveAlgebraic}($T$) \Comment{Alg.~\ref{alg:rr}; raw roots, no limits, may be empty}
  \State verified $\gets [\,]$
  \For{$q$ \textbf{in} candidates}
    \If{any entry of $q$ is not finite} \Continue{} \EndIf
    \State $r \gets \norm{\mathrm{FK}(q) - T}_F$ \Comment{Frobenius norm of the $4\times4$ difference}
    \If{$r \le 10^{-5}$} append $(q, r)$ to verified \EndIf \Comment{FK-closure gate; near-misses are dropped (no Newton polish)}
  \EndFor
  \State deduped $\gets [\,]$
  \For{$(q, r)$ \textbf{in} verified}
    \If{some $(q', r') \in$ deduped has $\max_i |\mathrm{wrap}(q_i - q'_i)| \le 10^{-3}$} keep the better one
    \Else{} append $(q, r)$
    \EndIf
  \EndFor \Comment{one representative per branch, $2\pi$ twins collapsed}
  \If{deduped $= \varnothing$ \textbf{and} \code{allow_rescue} \textbf{and} $\norm{p_T} \le$ reach radius}
    \State deduped $\gets$ \textsc{RescueViaPerturbation}($T$) \Comment{Alg.~\ref{alg:rescue}; reach radius = sum of link offsets}
  \EndIf
  \State sols $\gets$ \textsc{WrapToLimits}(deduped) \Comment{out-of-range joint: try $q_i + 2\pi k$, $k \in \{1,-1,2,-2\}$}
  \State sols $\gets$ \textsc{RespectLimits}(sols) \Comment{drop any solution with a joint still outside its range}
  \If{$q_{\text{seed}} \neq$ \textsc{None}} sort sols by $\max_i |\mathrm{wrap}(q_i - q_{\text{seed},i})|$ \EndIf \Comment{wrap-to-$\pi$ metric}
  \If{\code{max_solutions} given} truncate sols to the first \code{max_solutions} \EndIf
  \State \Return sols \Comment{each carries $q$, its FK residual, and a flag saying whether the rescue produced it}
\EndFunction
\end{algo}

\noindent
Three consequences of this ordering are worth holding in mind. (i) The limit
filter runs \emph{after} dedup, so a branch whose only in-limit representative
is a $2\pi$ twin still survives: \textsc{WrapToLimits} finds that twin. (ii)
\textsc{WrapToLimits} prefers the smallest $|k|$, i.e.\ it biases toward the
raw algebraic value, not toward the seed; the adapter's \textsc{Rebranch}
corrects that afterwards. (iii) With \code{max_solutions} set (the goal solve
in \code{_ik_dual_arm_at_frames} uses 8), the truncation happens under ssik's
wrap metric, \emph{before} the adapter re-sorts under the raw metric. Every
other call site passes \code{None} and receives the full branch set.

\begin{algo}{\textsc{SolveAlgebraic} --- Raghavan--Roth / Manocha--Canny core (\texttt{\_solve\_algebraic})}{alg:rr}
\Function{SolveAlgebraic}{$T$}
  \State $T_{\text{dh}} \gets T_{\text{pre}}^{-1}\, T\, T_{\text{post}}^{-1}$ \Comment{URDF (POE) frames $\to$ calibrated DH chain}
  \State $(P_{\sin}, P_{\cos}, P_{1}, Q) \gets$ \textsc{BuildPQ}($T_{\text{dh}}$) \Comment{$14\times9$ / $14\times8$; DH constants inlined at build time}
  \State $(E_{\sin}, E_{\cos}, E_{1}) \gets N^{\top} (P_{\sin}, P_{\cos}, P_{1})$ with $N^{\top} Q = 0$ \Comment{removes the $(q_1, q_2)$ terms; $6\times9$}
  \State $(E_{\text{quad}}, E_{\text{lin}}, E_{\text{const}}) \gets$ half-angle substitution $x_i = \tan(q_i/2)$ for $i = 3,4,5$
  \State $(M_{\text{quad}}, M_{\text{lin}}, M_{\text{const}}) \gets$ stack the six equations with $x_4$-shifted copies \Comment{$12\times12$ pencil in $x_3$}
  \State $(\text{roots}, \text{vecs}) \gets$ real eigenvalues of $M_{\text{quad}} x_3^2 + M_{\text{lin}} x_3 + M_{\text{const}}$ \Comment{M\"obius fallback if ill-conditioned}
  \State out $\gets [\,]$
  \For{$(x_3, v)$ \textbf{in} zip(roots, vecs)}
    \State read $q_4, q_5$ from the eigenvector $v$ \Comment{back-substitution starts}
    \State recover $q_1, q_2$ from the $P/Q$ system, then $q_6$
    \If{back-substitution failed} \Continue{} \EndIf
    \State append $q_{\text{dh}} - \theta_{\text{offset}}$ to out \Comment{back to the URDF joint zero}
  \EndFor
  \State \Return out \Comment{at most 16 candidates; each is re-checked by FK in Alg.~\ref{alg:artifact}}
\EndFunction
\end{algo}

\noindent
(Joint numbering in Alg.~\ref{alg:rr} is 1-based to match the text; the
source uses 0-based names such as \code{eliminate_q0_q1}.) The one-time cost
of choosing which joint to eliminate first was paid at \code{ssik build};
at run time the whole routine is dense linear algebra on small matrices.

\begin{algo}{\textsc{RescueViaPerturbation} --- \texttt{ssik.refinement.rescue} (only when \texttt{allow\_rescue=True})}{alg:rescue}
\Function{RescueViaPerturbation}{$T$}
  \State found $\gets [\,]$;\quad rng $\gets$ fixed seed \Comment{deterministic}
  \For{$i = 0 \dots 15$}
    \State $m \gets (1, 2, 4, 10)[i \bmod 4]$ \Comment{escalating scale}
    \State $\Delta T \gets$ random small motion: translation $\sim 5\,\text{mm} \cdot m$, rotation $\sim 5\,\text{mrad} \cdot m$
    \For{$q$ \textbf{in} \textsc{ArtifactSolve}($T\,\Delta T$; limits off, rescue off)} \Comment{off the ridge the algebra works}
      \State $q^{\star} \gets$ LM-refine $q$ toward $T$ (analytic Jacobian, $\le 20$ iterations)
      \If{$q^{\star}$ failed \textbf{or} $\norm{\mathrm{FK}(q^{\star}) - T} > 10^{-8}$} \Continue{} \EndIf
      \If{$q^{\star}$ within $10^{-4}$ rad (wrap-to-$\pi$) of something already in found} \Continue{} \EndIf
      \State append $q^{\star}$ to found, tagged \code{refinement_used="lm"}
    \EndFor
  \EndFor
  \State \Return found
\EndFunction
\end{algo}

\noindent
The rescue exists for a measure-zero situation: a reachable pose at which the
Raghavan--Roth pencil loses rank, so the algebra extracts nothing even though
solutions exist arbitrarily close by. It costs about sixteen ordinary solves
plus refinement, roughly $40\times$ a normal call according to the loader's
own comment. That is why the planner passes \code{allow_rescue=False} inside
every per-waypoint loop (tree extension, tracking, stitching, smoothing) and
leaves it on only for the one-off endpoint solves: in a loop an unreachable
waypoint should fail in about a millisecond, and a rescued, numerically
refined solution may sit on a different branch than its neighbours anyway.

\paragraph{What an empty list means.}
\code{ssik_arm_branches} returns $[\,]$ in exactly three cases: the tool pose
is outside the arm's reach (no real roots close under FK), every branch has a
joint outside its limits even after $\pm 2\pi$ wrapping, or -- with rescue off
-- the pose sits on a rank-deficient ridge. The planner does not distinguish
these; all three become \code{ik_failure} in the profile counters.

\paragraph{Where each call site sits.}
\begin{center}
\footnotesize
\begin{tabular}{@{}l l l c l@{}}
\toprule
caller & used for & seed & \code{max_solutions} & rescue \\
\midrule
\code{_solve_dual_arm_pose_ik_ssik} & every waypoint & previous $q$ & all & off \\
\code{_ik_dual_arm_at_frames} & goal from tool frames & placeholder start & 8 & on \\
\code{_ssik_pair_goal_branch_with_home} & branch pairing & none & all & on \\
\code{_solve_endpoint_dual_arm_ik_ssik} & goal repair (Alg.~\ref{alg:endpoint}) & goal $q$ & 8 & on \\
\bottomrule
\end{tabular}
\end{center}

\subsection{Growing one tree toward a pose}

Algorithm~\ref{alg:extend} is used both for the ``extend toward a random
sample'' step and for the ``connect the other tree'' step. It always walks the
whole segment; it does not take a single $\varepsilon$ step.

\begin{algo}{\textsc{ExtendToward} --- \texttt{extend\_toward}}{alg:extend}
\Function{ExtendToward}{$\Tree$, $v_{\text{src}}$, $T_{\text{target}}$}
  \State $v \gets v_{\text{src}}$;\quad $q \gets \Tree.\text{conf}[v]$;\quad reached $\gets$ \textbf{true};\quad reason $\gets$ ``reached''
  \State \textbf{define} \fn{Stop}($r$): reason $\gets r$;\ reached $\gets$ \textbf{false};\ leave the loop
  \For{\D{$T$ \textbf{in} \fn{Interp}($v.\text{pose}$, $T_{\text{target}}$; $\delta_p$, $\delta_r$)}} \dtag{D-5}
    \If{\D{IK enabled} (stage $\ge 2$)} \dtag{D-6}
      \State \D{$q' \gets$ \textsc{SolveDualArmPoseIK}($T$, $q$)} \Comment{warm-seeded from the parent}
      \If{\D{$q' = $ \textsc{None}}} \fn{Stop}(``ik\_failure'') \EndIf
      \If{\D{\fn{Jump}($q'$, $q$)}} \fn{Stop}(``continuity'') \dtag{D-6} \EndIf
      \If{\D{stage $= 3$ \textbf{and} \code{joint_collision_fn}($q'$)}} \fn{Stop}(``collision'') \EndIf
    \EndIf
    \If{\D{\code{collision_fn}($T$)}} \fn{Stop}(``collision'') \Comment{bar-only check} \EndIf
    \State \D{$v_{\text{new}} \gets$ new node(pose $T$, parent $v$); add to $\Tree$; cache conf $q'$ and $\phi(T)$} \dtag{D-1}
    \State $v \gets v_{\text{new}}$;\quad $q \gets q'$
  \EndFor
  \State \Return ($v$, reached, reason) \Comment{$v$ is the last node that was accepted}
\EndFunction
\end{algo}

\noindent
Because the loop appends every accepted waypoint before it fails, a blocked
walk still leaves a chain of nodes reaching to the obstacle. That is the same
behaviour as the baseline's \fn{Connect} (which calls \fn{Extend} repeatedly),
but here it is also what the \emph{first} extension does.

\subsection{The bidirectional loop}

\begin{algo}{\textsc{PlanPoseBiRRT} --- \texttt{plan\_pose\_birrt}}{alg:birrt}
\Function{PlanPoseBiRRT}{$T_s$, $T_g$, $q_s$, $q_g$}
  \State \D{\textbf{require} $q_s$, $q_g$ hold the bar at $T_s$, $T_g$ and pass \code{joint_collision_fn}} \dtag{D-9}
  \Statex \hspace{\algorithmicindent}\textcolor{gray}{\footnotesize (stage 1 checks only \code{collision_fn}($T_s$), \code{collision_fn}($T_g$))}
  \For{attempt $= 1 \dots$ \code{max_attempts}} \dtag{D-8}
    \State \D{$\Tree_a \gets$ tree rooted at ($T_s$, $q_s$, $\phi(T_s)$);\quad $\Tree_b \gets$ tree rooted at ($T_g$, $q_g$, $\phi(T_g)$)} \dtag{D-1}
    \For{iter $= 1 \dots$ \code{max_iterations}}
      \If{elapsed $\ge$ \code{max_time}} \Break \EndIf
      \State \D{$(\Tree_{\text{grow}}, \Tree_{\text{other}}) \gets (\Tree_a, \Tree_b)$ \textbf{if} $|\Tree_a| \le |\Tree_b|$ \textbf{else} $(\Tree_b, \Tree_a)$} \dtag{D-4}
      \If{\D{rand() $<$ \code{goal_sample_prob} ($=0.05$)}} \dtag{D-3}
        \State \D{$T_{\text{rand}} \gets$ root pose of $\Tree_{\text{other}}$} \Comment{cross-bias toward the other tree}
      \Else
        \State \D{$T_{\text{rand}} \gets$ \textsc{SamplePose}()} \Comment{uniform box around the base, random Euler angles}
      \EndIf
      \State \D{$v_{\text{near}} \gets \arg\min_{v \in \Tree_{\text{grow}}} \norm{\phi(v.\text{pose}) - \phi(T_{\text{rand}})}_2$} \dtag{D-2}
      \State \D{$(v_{\text{new}}, \_, \text{reason}) \gets$ \textsc{ExtendToward}($\Tree_{\text{grow}}$, $v_{\text{near}}$, $T_{\text{rand}}$)} \dtag{D-5}
      \State record reason in \code{extend_stop_reasons}
      \If{$v_{\text{new}} = v_{\text{near}}$} \Continue{} \EndIf \Comment{no progress at all}
      \State \D{$u_{\text{near}} \gets \arg\min_{u \in \Tree_{\text{other}}} \norm{\phi(u.\text{pose}) - \phi(v_{\text{new}}.\text{pose})}_2$}
      \State \D{$(u_{\text{new}}, \text{connected}, \_) \gets$ \textsc{ExtendToward}($\Tree_{\text{other}}$, $u_{\text{near}}$, $v_{\text{new}}.\text{pose}$)} \dtag{D-5}
      \If{\textbf{not} connected} \Continue{} \EndIf
      \Statex
      \Statex \hspace{\algorithmicindent}\hspace{\algorithmicindent}\textcolor{gray}{\footnotesize --- the two trees now share the pose $v_{\text{new}}.\text{pose}$; order the path start $\to$ goal ---}
      \If{$\Tree_{\text{grow}} = \Tree_a$}
        \State $S \gets$ \fn{Retrace}($v_{\text{new}}$) \Comment{root of $\Tree_a$ (start) $\to$ seam}
        \State $E \gets$ \fn{Reverse}(\fn{Retrace}($u_{\text{new}}$)) \Comment{seam $\to$ root of $\Tree_b$ (goal)}
      \Else
        \State $S \gets$ \fn{Retrace}($u_{\text{new}}$);\quad $E \gets$ \fn{Reverse}(\fn{Retrace}($v_{\text{new}}$))
      \EndIf
      \State $E \gets E[1{:}]$ \Comment{drop the duplicated seam node}
      \State $\text{poses} \gets$ poses of $S$ followed by poses of $E$
      \If{IK enabled}
        \State \D{$\text{confs} \gets$ \textsc{Stitch}($S$, $E$, $q_g$)} \dtag{D-7}
        \If{\D{confs $=$ \textsc{None}}} \Continue{} \EndIf \Comment{seam not joint-continuous; keep searching}
      \EndIf
      \State \Return (poses, confs)
    \EndFor
    \State \Comment{attempt exhausted: discard both trees and start fresh}
  \EndFor
  \State \Return (\textsc{None}, \textsc{None})
\EndFunction
\end{algo}

\subsection{Stitching the seam}

Each tree carries its own joint configurations, grown outward from its root by
warm-seeded IK. The start tree's configurations therefore sit on the branch
sheet of $q_s$ and the goal tree's on the sheet of $q_g$. Even when the two
trees meet at the same bar pose, the two configurations at the seam may
differ by a branch flip. Algorithm~\ref{alg:stitch} repairs this by re-solving
one side from the other side's seam configuration.

\begin{algo}{\textsc{Stitch} --- inner functions \texttt{\_stitch\_forward} / \texttt{\_stitch\_backward}}{alg:stitch}
\Function{Stitch}{$S$, $E$, $q_g$}
  \State $Q_S \gets$ cached confs of $S$;\quad $Q_E \gets$ cached confs of $E$
  \State \D{$Q_E' \gets$ \textsc{Rewalk}(poses of $E$, seed $= Q_S[-1]$)} \Comment{forward: goal side on the start sheet}
  \If{$Q_E' \neq$ \textsc{None}} \Return $Q_S \,\Vert\, Q_E'$ \EndIf
  \State \D{$Q_S' \gets$ \textsc{Reverse}(\textsc{Rewalk}(\textsc{Reverse}(poses of $S$), seed $= Q_E[-1]$))}
  \Statex \hspace{\algorithmicindent}\textcolor{gray}{\footnotesize backward: start side rebuilt on the goal sheet (the seed falls back to $q_g$ when $E$ is empty, i.e.\ the trees met at the goal root)}
  \If{$Q_S' = $ \textsc{None} \textbf{or} \fn{Jump}($Q_S'[0]$, $Q_S[0]$)} \Return \textsc{None} \EndIf \Comment{must still end at the given start conf}
  \State \Return $Q_S' \,\Vert\, Q_E$
\EndFunction
\Statex
\Function{Rewalk}{poses, seed}
  \State $q \gets$ seed;\quad out $\gets [\,]$
  \For{$T$ \textbf{in} poses}
    \State $q' \gets$ \textsc{SolveDualArmPoseIK}($T$, $q$)
    \If{$q' = $ \textsc{None} \textbf{or} \fn{Jump}($q'$, $q$) \textbf{or} \code{joint_collision_fn}($q'$)} \Return \textsc{None} \EndIf
    \State append $q'$ to out;\quad $q \gets q'$
  \EndFor
  \State \Return out
\EndFunction
\end{algo}

\noindent
Note what a stitch failure does \emph{not} do: it does not remove any node
from either tree. The iteration simply continues, and the same seam can be hit
again later from a different frontier. The failure is counted under
\code{stitch_ik_failure}, \code{stitch_continuity_fail},
\code{stitch_collision} or \code{stitch_endpoint_mismatch} in the profile.

\subsection{Sampling}

\textsc{SamplePose} (\code{sample_pose} with \code{goal_sample_prob}$=0$ in
the bidirectional caller) draws the bar position uniformly in a box centred on
the robot base: $\pm 1.1$ m in $x$ and $y$ (\code{workspace_xy}$=2.2$), and
$z \in [\max(0.05, z_{\text{base}}),\ \max(0.05, z_{\text{base}}) + 1.2]$.
Orientation is built from three independent uniform angles in $[-\pi, \pi]$
(roll, pitch, yaw). This is \emph{not} uniform on the rotation group, but no
part of the planner depends on it being so. Note that there is no bias toward
the goal pose in the random branch; the only bias is the 5\,\% cross-bias
toward the \emph{other tree's root} in Algorithm~\ref{alg:birrt}.

% =============================================================================
\section{Core differences from the standard BiRRT}
\label{sec:diff}
% =============================================================================

The tags below match the red \D{\textsf{[D-n]}} markers in the pseudocode.

\begin{description}[leftmargin=1.6cm,style=nextline,itemsep=4pt]
  \item[D-1 \ Nodes are bar poses, not configurations.]
    A tree node stores a pose $T \in \SE$ plus two side caches: the joint
    configuration that IK produced for it, and its feature vector
    $\phi(T)$. The search never samples joints. The joint path is a
    \emph{by-product} that is grown alongside the pose tree, one warm-seeded IK
    per node.

  \item[D-2 \ Distance is measured on the bar's corners.]
    Nearest-neighbour search uses $\norm{\phi(\cdot)-\phi(\cdot)}_2$
    (Eq.~\ref{eq:feature}), a Euclidean distance on eight projected corner
    points. Translation and rotation are combined through the geometry of the
    carried object instead of a hand-chosen weight.

  \item[D-3 \ No goal bias; a small cross-bias instead.]
    Random samples are drawn from a fixed box around the base. With
    probability $0.05$ the target is replaced by the \emph{other} tree's
    root, which pulls the two trees toward each other. The single-tree
    planner (\code{plan_pose_rrt}, archival) used a $0.1$ goal bias.

  \item[D-4 \ The smaller tree grows.]
    Instead of strict alternation, each iteration extends whichever tree
    currently has fewer nodes, then asks the other tree to connect. The path
    is always assembled start $\to$ goal regardless of which tree grew.

  \item[D-5 \ Extend is a full greedy walk with partial progress kept.]
    Both the extension toward the sample and the connection toward the new
    frontier walk the entire interpolated segment at resolution
    ($\delta_p$, $\delta_r$), stopping at the first rejected waypoint and
    keeping every accepted one. The baseline takes one fixed step for the
    first move. ``Connected'' therefore means the other tree walked all the
    way to the exact pose of the new frontier (the last interpolated pose is
    the target itself), not that it came within a tolerance.

  \item[D-6 \ Four validity gates per waypoint, one of them kinematic.]
    A waypoint is accepted only if (a) both arms have an analytical IK
    solution within joint limits, (b) the nearest solution is within $\tau =
    10^\circ$ per joint of the previous waypoint's configuration, (c) that
    configuration is collision-free for the whole cell (stage 3), and (d) the
    bar alone is collision-free at that pose. Gate (b) has no counterpart in
    the baseline; it is what turns a pose path into a joint path that can be
    executed without a branch flip, and it is the dominant reason extensions
    stop (Section~\ref{sec:bench}).

  \item[D-7 \ Meeting at a pose is not enough; the seam is re-solved.]
    The two trees carry joint configurations from two different roots, so
    when they meet the joint values may disagree even though the poses are
    identical. The goal-side configurations are rebuilt from the start-side
    seam configuration (or, failing that, the start side from the goal
    branch). If neither rewalk passes the continuity and collision gates, the
    candidate path is thrown away and the search continues with both trees
    intact.

  \item[D-8 \ Restarts with fresh trees.]
    Up to \code{max_attempts} independent attempts, each with its own time
    budget (\code{max_time}) and iteration cap. Trees are not carried over
    between attempts; the random generator is.

  \item[D-9 \ Both endpoints must be given as configurations.]
    The goal tree needs a root configuration, so $q_g$ has to be solved before
    the search. Because the whole path must stay on one branch sheet, the
    endpoint IK cannot be chosen independently for start and goal; the
    pre-processing in Section~\ref{sec:endpoints} exists to make $q_s$ and
    $q_g$ mutually reachable.
\end{description}

% =============================================================================
\section{Around the search: endpoints, stages, smoothing}
\label{sec:endpoints}
% =============================================================================

\subsection{Stages}

\code{plan_constrained_dual_arm(stage=...)} switches gates on incrementally,
which is how the planner was built and how it is still debugged:
\begin{center}
\small
\begin{tabular}{@{}clll@{}}
\toprule
stage & IK per waypoint & continuity gate & collision \\
\midrule
1 & off & off & bar only (\code{collision_fn}) \\
2 & on  & on  & bar only \\
3 & on  & on  & bar + full cell (\code{joint_collision_fn}) \\
\bottomrule
\end{tabular}
\end{center}
Stage 3 is the production setting. A stage-1 success returns
\code{info["pose_only_success"]} and no joint path.

\subsection[Getting the two endpoints (derive\_start=True)]{Getting the two endpoints (\texttt{derive\_start=True})}

For the M1 movement the Rhino export carries the goal as two tool0 frames and
leaves the start joints as a placeholder, so the entry point computes both
endpoints itself. The order is \emph{goal first}: the goal fixes the grasp
transforms, then a start is found that the goal can reach.

\begin{enumerate}[nosep]
  \item \textbf{Goal IK.} Solve each arm at its goal tool frame, nearest
        analytical branch to the (placeholder) seed; FK at that configuration
        gives $T_g$ and, with the authored attachment frame, $G_L$ and $G_R$.
  \item \textbf{Goal/home branch pairing}
        (\code{_ssik_pair_goal_branch_with_home}). Enumerate all legal IK
        branches at the goal and at every candidate home pose (three carry
        anchors $\times$ their roll/yaw variants), rank each goal branch by
        its $\norm{\cdot}_\infty$ distance to the nearest home branch after
        $2\pi$ re-branching, and take the first of the top-$4 \times 4$ pairs
        that is FK-valid and collision-free. Measured on B6, independently
        picked branches were $104$--$230^\circ$ apart; the best pairing is
        $38^\circ / 67^\circ$.
  \item \textbf{Goal collision repair.} If the paired goal collides,
        \code{solve_endpoint_dual_arm_ik} is asked for a different
        configuration at the \emph{same} bar pose and grasps. Its signature
        still carries \code{rng} and \code{max_attempts} from the gradient
        era, where PyBullet's numeric IK converges to whichever branch basin
        its seed lands in, so re-seeding at random was the only way to reach
        another branch. On the ssik backend both arguments are ignored: the
        branch set is finite and enumerable, so the ssik variant scans it
        deterministically (Alg.~\ref{alg:endpoint}). ``Re-solve'' therefore
        means ``walk the list of left$\times$right branch pairs, nearest to
        the current goal first, and take the first one that holds the bar and
        is collision-free''. If none does, the plan fails with
        \code{goal_in_collision}.
  \item \textbf{Backward-tracked start} (Alg.~\ref{alg:tracked}): walk the
        bar from $T_g$ toward candidate home poses with the same warm-seeded IK
        and continuity gate the tree uses. Whatever it arrives at is on the
        goal's own branch sheet by construction.
  \item \textbf{Direct-path shortcut.} If one of those walks is also
        collision-free at every waypoint, its reversal \emph{is} the M1 path
        and the search is skipped entirely; the returned
        \code{info["planner"]} then reads \code{"direct_path"}. The smoother
        is not run in that case.
  \item \textbf{Cold fallback.} \code{derive_constrained_start}, the older
        endpoint sweep, is used when tracking finds nothing: it picks a
        holdable home pose by endpoint IK alone.
\end{enumerate}

\begin{algo}{\textsc{EndpointIK} --- \texttt{\_solve\_endpoint\_dual\_arm\_ik\_ssik} (goal collision repair)}{alg:endpoint}
\Function{EndpointIK}{$T$, $G_L$, $G_R$, $q_{\text{seed}}$, \code{collision_fn}}
  \State $B_L \gets$ \textsc{SsikArmBranches}(L, $T G_L$, $q_{\text{seed},L}$, \code{max_solutions}$=8$, rescue on)
  \State $B_R \gets$ \textsc{SsikArmBranches}(R, $T G_R$, $q_{\text{seed},R}$, \code{max_solutions}$=8$, rescue on)
  \If{$B_L = \varnothing$ \textbf{or} $B_R = \varnothing$} \Return \textsc{None} \EndIf \Comment{an arm cannot reach its grasp at all}
  \State pairs $\gets [\,(q_L, q_R) \text{ for } q_L \in B_L,\ q_R \in B_R\,]$ \Comment{up to $8 \times 8 = 64$ combinations}
  \State sort pairs by $\norm{(q_L, q_R) - q_{\text{seed}}}_2$ ascending \Comment{stay close to the goal the pairing chose}
  \For{$(q_L, q_R)$ \textbf{in} pairs}
    \State $q \gets (q_L, q_R)$
    \If{\textbf{not} \code{validate_dual_arm_bar_pose}($q$, $T$)} \Continue{} \EndIf \Comment{PyBullet FK gate}
    \If{\code{collision_fn}($q$)} \Continue{} \EndIf \Comment{full-cell check}
    \State \Return $q$ \Comment{first pair that is reachable, consistent, and free}
  \EndFor
  \State \Return \textsc{None} \Comment{every branch pair is either FK-inconsistent or colliding}
\EndFunction
\end{algo}

\noindent
Two details. The pair order uses the Euclidean norm on the raw 12-vector,
not the $\norm{\cdot}_\infty$ of the continuity gate, and nothing here checks
continuity: the repaired goal may sit on a different branch than the one
\code{_ssik_pair_goal_branch_with_home} chose, in which case the
backward-tracked start derivation (next step) starts from the repaired goal
and finds a home on \emph{its} sheet. With the gradient backend the same
function instead loops \code{max_attempts} times over
\code{solve_dual_arm_pose_ik}, seeding attempt 0 from $q_{\text{seed}}$ and
every later attempt from a uniform random 12-vector in $[-\pi, \pi]$.

\begin{algo}{\textsc{DeriveStartTracked} --- \texttt{derive\_constrained\_start\_tracked} (condensed)}{alg:tracked}
\Function{DeriveStartTracked}{$T_g$, $q_g$}
  \State variants $\gets$ every selected anchor's canonical orientation, then roll/yaw variants smallest-first
  \State deltas $\gets$ grid in the sweep box ($\pm0.3$ m, step $0.1$ m), sorted nearest-first
  \State budget $\gets 120$ s split evenly across the selected anchors;\quad first\_start, best\_partial $\gets$ \textsc{None}
  \For{(label, $T_{\text{anchor}}$) \textbf{in} variants}
    \For{$\Delta$ \textbf{in} (all deltas \textbf{if} canonical \textbf{else} first 60)}
      \If{this anchor's time share is spent} \Break \EndIf
      \State $T_h \gets T_{\text{anchor}}$ shifted by $\Delta$
      \State (track, ok) $\gets$ \textsc{TrackDirectPath}($T_g$, $q_g$, $T_h$; screening step $0.01$ m / $0.025$ rad)
      \Statex \hspace{\algorithmicindent}\hspace{\algorithmicindent}\hspace{\algorithmicindent}
        \textcolor{gray}{\footnotesize = \textsc{Rewalk} without the collision gate; returns the accepted prefix and whether it reached $T_h$}
      \If{\textbf{not} ok}
        \State best\_partial $\gets$ the broken track whose last pose is farthest from $T_g$ \Continue{}
      \EndIf
      \If{\code{joint_collision_fn}(track[-1].conf)} \Continue{} \EndIf \Comment{the arrival is a plan endpoint}
      \If{no interior waypoint of track collides (scanned every 3rd first)}
        \State re-track at the requested ($\delta_p$, $\delta_r$) if finer than screening; on success
        \State \Return ($T_h$, arrival conf, \code{direct_path} = track) \Comment{caller skips the search}
      \EndIf
      \If{first\_start $=$ \textsc{None}} first\_start $\gets$ ($T_h$, arrival conf) \EndIf \Comment{blocked mid-way; keep as start-only}
    \EndFor
  \EndFor
  \If{first\_start $\neq$ \textsc{None}} \Return first\_start \EndIf
  \If{best\_partial reaches $\ge 2$ cm from $T_g$} \Comment{M0 will cover the remaining gap}
    \State \Return its last collision-free waypoint still $\ge 2$ cm from $T_g$, flagged \code{partial}
  \EndIf
  \State \Return (\textsc{None}, \textsc{None})
\EndFunction
\end{algo}

\subsection{Smoothing}

The raw path is a chain of interpolated waypoints with sharp corners at every
tree node. The smoother (Algorithm~\ref{alg:smooth}) is a classic random
shortcut in \emph{pose} space with the joint path re-derived for each accepted
shortcut, so the same three gates (IK, continuity, collision) hold on the
result.

\begin{algo}{\textsc{SmoothPosePath} --- \texttt{smooth\_dual\_arm\_pose\_path}}{alg:smooth}
\Function{SmoothPosePath}{poses, confs}
  \State cost $\gets \sum_i d(T_i, T_{i+1})$ (Eq.~\ref{eq:feature});\quad corners $\gets$ indices where the feature-vector direction changes
  \For{iter $= 1 \dots$ \code{smooth_max_iterations} ($=100$), within \code{smooth_max_time} ($=10$ s)}
    \If{fewer than 3 corners} \Break \EndIf
    \State pick two corners $i < j$ at least two corners apart
    \State shortcut $\gets$ \fn{Interp}($T_i$, $T_j$; $\delta_p$, $\delta_r$);\quad candidate $\gets$ poses$[{:}i] \,\Vert\,$ shortcut $\Vert$ poses$[j{+}1{:}]$
    \If{cost(candidate) $\ge$ cost} \Continue{} \EndIf
    \State $Q \gets$ \textsc{Rewalk}(candidate$[i{:}]$, seed $=$ confs$[i]$) \Comment{IK + continuity + collision from $i$ to the end}
    \If{$Q = $ \textsc{None}} \Continue{} \EndIf
    \State poses $\gets$ candidate;\quad confs $\gets$ confs$[{:}i] \,\Vert\, Q$;\quad recompute cost and corners
  \EndFor
  \State \Return (poses, confs)
\EndFunction
\end{algo}

\noindent
The rewalk starts at $i$ and runs to the \emph{end} of the path, not only over
the shortcut, because the configurations after $j$ were seeded from the old
waypoints and may no longer be continuous with the new ones.

% =============================================================================
\section{Parameters and their defaults}
\label{sec:params}
% =============================================================================

The same knob often has a different default at each layer. The values below
are what the code carries today.

\begin{center}
\footnotesize
\setlength{\tabcolsep}{4pt}
\begin{tabular}{@{}p{3.6cm} p{2.6cm} p{2.6cm} p{2.8cm} p{3.8cm}@{}}
\toprule
parameter & core default & api default & headless M1 & meaning \\
\midrule
\code{position_res} $\delta_p$ & 0.05 m & 0.01 m & 0.002 m & extend step; also the goal tolerance \\
\code{rotation_res} $\delta_r$ & 0.10 rad & 0.025 rad & 0.002 rad & extend step; also the goal tolerance \\
\code{max_time} & 30 s & 30 s & CLI \code{--max-time} & per attempt \\
\code{max_iterations} & 2000 & 2000 & --- & per attempt \\
\code{max_attempts} & 5 & 5 & --- & fresh tree pairs \\
\code{goal_sample_prob} & 0.05 & (not exposed) & --- & cross-bias to the other root \\
\code{workspace_xy} / \code{_z} & 2.2 m / 1.2 m & (not exposed) & --- & sampling box \\
$\tau$ (continuity) & $10^\circ$ & $10^\circ$ & --- & the constant in Eq.~\ref{eq:gate} \\
\code{dist_metric} & \code{"feature"} & (not exposed) & --- & Eq.~\ref{eq:feature} \\
smoothing & --- & 100 iters / 10 s & --- & smoother iteration and time caps \\
tracked-start budget & 120 s & (not exposed) & --- & split evenly across anchors \\
screening step & 0.01 m / 0.025 rad & (not exposed) & --- & never finer than $\delta_p,\delta_r$ \\
\code{start_home_anchor} & --- & all & \code{--cdfm-home-anchor} & horizontal / vertical / back \\
\bottomrule
\end{tabular}
\end{center}

\noindent
Here \emph{core} is \code{plan_pose_birrt} (and, for the last three rows,
\code{derive_constrained_start_tracked}), \emph{api} is
\code{plan_constrained_dual_arm}, and \emph{headless M1} is what
\path{scripts/headless_bar_action_planner.py} passes for an M1 movement.
\medskip

\noindent
$\delta_p$ and $\delta_r$ reach four places: the tree extension, the connect
walk, the smoother's shortcuts, and the delivered direct path. Lowering them
makes every one of those denser; it also makes each extension cost more IK
solves and collision checks.

% =============================================================================
\section{What the benchmark measured}
\label{sec:bench}
% =============================================================================

These numbers are copied from \code{benchmark/analysis_night1.md} (run
\code{night1}: 24{,}000 runs over 16 bars, 800 goals, 3 starts per goal, 10
seeds; each run is a coarse pass at $0.1$ m / $0.1$ rad, 30 s and a fine pass
at $0.002$ m / $0.002$ rad, 300 s). They are observations, not conclusions,
and that document states nothing in it has been acted on.

\begin{itemize}[nosep]
  \item Solved: 16.3\,\% of runs; 28.0\,\% of problems within ten seeds.
        Median solved time 143 s (fine), 21 s (coarse).
  \item Why extensions stop: \textbf{continuity 75.4\,\%}, IK failure
        18.5\,\%, collision 3.7\,\%. The gate in D-6(b) is the main brake,
        not the environment.
  \item Where the time goes: collision checking 36.8\,\% (about 12.9 ms per
        call), IK 33.0\,\%, stitching 13.0\,\%, smoothing 9.8\,\%.
  \item The tracked start derivation found a collision-free direct path (and
        skipped the search) for 34.9\,\% of goals and used its whole 120 s
        budget at the median; the survey found that 73\,\% of goals have at
        least one straight-line-solvable start.
  \item On problems where a straight-line motion provably exists, the search
        still failed 56\,\% of the time, for the same continuity reason.
\end{itemize}

\paragraph{Where the ``why extensions stop'' numbers come from.}
The three percentages are lifted from Finding~1 of
\path{benchmark/analysis_night1.md}; the raw run data behind them
(\path{runs/night1/summary.json} under the benchmark data root) is not in
the repository and was not available on the machine this report was written
on, so the shares could not be recomputed here. What \emph{can} be traced is
the counter itself:
\begin{enumerate}[nosep]
  \item \code{plan_pose_birrt} keeps one \code{Counter} named
        \code{extend_stop_reasons}. It is incremented \textbf{once per call
        to \textsc{ExtendToward}}, with the walk's stop reason:
        \code{reached}, \code{collision}, \code{ik_failure} or
        \code{continuity} for the growth walk (Alg.~\ref{alg:birrt}
        line~17), the same names prefixed \code{connect_} for the connect
        walk, and the four keys
        \texttt{stitch\_\{ik\_failure, continuity\_fail, collision, endpoint\_mismatch\}}
        from Alg.~\ref{alg:stitch}.
  \item \code{update_debug_tree} copies the counter into the caller's
        \code{debug_tree_out} dict, and the entry point returns that dict as
        \code{info["profile"]}. The histogram is the entry\newline
        \code{"extend_stop_reasons"} of that dict.
  \item The benchmark (\path{benchmark/cdfm.py}) stores that dict on every
        pass record, and \path{benchmark/summarize.py} adds the dicts of
        \emph{both} passes (coarse and fine) of \emph{every} run (solved or
        not) into one histogram per group.
\end{enumerate}
Two consequences for reading the table. A ``continuity'' entry is one
\emph{walk} that ended on the gate, not one waypoint: a walk that added twenty
nodes and then hit a branch flip counts exactly like one that died on its
first step. And the analysis document does not say which keys form its
denominator; the dashboard plots every key as-is, so \code{reached},
\code{connect_*} and \code{stitch_*} are presumably what ``everything else,
2.4\,\%'' contains. If so, connect walks are rare, which would mean most
growth walks make no progress at all before the gate fires. That reading
should be confirmed against the summary file before anything is built on it.

\paragraph{Reading the numbers against the pseudocode.}
Every one of the three gates in Algorithm~\ref{alg:extend} is evaluated
\emph{per waypoint}, so at $\delta = 0.002$ a single extension across the
workspace is hundreds of IK solves and collision checks, and a single
rejected waypoint ends it. The continuity share suggests that ``nearest
branch to the seed'' alone does not keep the walk on one sheet often enough;
whether that is a threshold question, a sampling question, or an IK question
is exactly what the stage switches and the profile counters are there to
separate.

% =============================================================================
\section{Reading guide}
\label{sec:reading}
% =============================================================================

Suggested order for reading the code with this report open. Line numbers are
from the tree at commit \code{cd239d5} (29 September 2026).

\begin{center}
\footnotesize
\begin{tabular}{@{}p{5.8cm} p{6.6cm} p{3.4cm}@{}}
\toprule
read & at & this report \\
\midrule
\code{joint_step_exceeds_threshold} & \code{core.py:186} & Eq.~\ref{eq:gate} \\
\code{get_bar_feature_points}, \code{pose_to_feature_vec} & \code{core.py:212}, \code{:412} & Eq.~\ref{eq:feature} \\
\code{sample_pose}, \code{nearest_node} & \code{core.py:541}, \code{:583} & D-2, D-3 \\
\code{build_cfab_pose_collision_fn} & \code{core.py:645} & bar-only predicate \\
\code{validate_dual_arm_bar_pose} & \code{core.py:905} & FK gate in Alg.~\ref{alg:ik} \\
\code{_solve_dual_arm_pose_ik_ssik} & \code{core.py:1055} & Alg.~\ref{alg:ik} \\
\code{ssik_arm_branches}, \code{rebranch_toward_seed} & \code{ssik_ik.py:145}, \code{:118} & Alg.~\ref{alg:branches} \\
\code{ssik_inprocess.solve} & \code{keyframe/ssik_inprocess.py:99} & loader \\
artifact \code{solve}, \code{_solve_algebraic} & \code{asset/ssik/left_ur_arm_ik.py} \code{:1066}, \code{:817} & Algs.~\ref{alg:artifact}, \ref{alg:rr} \\
\code{wrap_to_limits}, \code{nearest_to_seed} & \code{ssik/postprocess.py:99}, \code{:157} & limit and seed filters \\
\code{rescue_via_T_perturbation} & \code{ssik/refinement/rescue.py:57} & Alg.~\ref{alg:rescue} \\
\code{extend_toward} & \code{core.py:2285} & Alg.~\ref{alg:extend} \\
\code{plan_pose_birrt} & \code{core.py:2775} & Alg.~\ref{alg:birrt}, \ref{alg:stitch} \\
\code{_solve_endpoint_dual_arm_ik_ssik} & \code{core.py:1292} & Alg.~\ref{alg:endpoint} \\
\code{track_direct_path}, \code{first_blocked_index} & \code{core.py:1824}, \code{:1905} & Alg.~\ref{alg:tracked} \\
\code{derive_constrained_start_tracked} & \code{core.py:1936} & Alg.~\ref{alg:tracked} \\
\code{reconstruct_joint_path_for_pose_path} & \code{core.py:2448} & \textsc{Rewalk} \\
\code{update_debug_tree} & \code{core.py:2511} & stop-reason export (Section~\ref{sec:bench}) \\
\code{smooth_dual_arm_pose_path} & \code{smooth.py:20} & Alg.~\ref{alg:smooth} \\
\code{_ssik_pair_goal_branch_with_home} & \code{api.py:254} & Section~\ref{sec:endpoints} step 2 \\
\code{_derive_constrained_start_for_plan} & \code{api.py:383} & Section~\ref{sec:endpoints} \\
\code{plan_constrained_dual_arm} & \code{api.py:753} & entry point \\
\bottomrule
\end{tabular}
\end{center}

\noindent
Two things the reader should not look for, because they are absent by design:
there is no goal-region test in the bidirectional planner (the trees meet at
an exact pose, D-5), and there is no node pruning or rewiring of any kind
(a failed stitch leaves both trees untouched, D-7).

\end{document}
